\section[第5章\q 踏出前的质疑]{第5章\\ 踏出前的质疑}
\label{sec:ch5}

这天柔伊在工作时，忍不住回想之前和亨利的谈话。后来，她拿出手机滑动屏幕。通勤到公司时，她在手机上打了一些文本，现在她重新誊写到电脑上，分门别类，并做了一些修饰。

这就是柔伊的工作日常。她是一位编辑，她有本事把这些零碎的文句，穿针引线重新编排，就像考古学家，可以从出土骨骼中，拼凑出一个远古时代。下标、修辞，那是下一步才要考虑的。第一步要考虑的是愿景：作者想说的重点是什么？

{\kaishu 大多数的人，赚得多的时候，花得更多。}

{\kaishu 每个人都觉得听过这话，但实践的人却很少。}

{\kaishu 把自己放在第一位。}

{\kaishu 有人说，这是宇宙间最强大的力量。}

{\kaishu 1天10美元，人生大不同。}

{\kaishu 让每天第一个小时的收入超过所值}

还有昨天记下来的话，很像一个理财大师的格言：

{\kaishu 如果你可以负担得起拿铁，你就可以负担得起这张照片。}

现在对这句话，或者说拿铁因子的说法，柔伊依然并不十分明白。

她的拿铁因子？

一天结束后，下了L线走回家时，她在海伦娜咖啡厅停了下来，她想看看亨利在不在，可以再向他请教一下拿铁因子的事。但是他已经下班了。他是下午3:00时离开咖啡厅的，柜台后的年轻店员告诉柔伊。

“我应该想到的。”柔伊说，毕竟他早上7:00就来了。“我猜想他是上早班吧！”

“早班？”年轻店员笑说：“亨利不轮班的。”

不轮班？“所以，通常他几点下班？”柔伊问。

“他想哪时候走就哪时候走。”年轻人说，耸了耸肩，“通常是3:00，但是也可能更晚一点，也可能更早一点。”

他哪时候想走就哪时候走？这是什么工作？

柔伊回到家，打开前门时，脑中依然在想着这个问题。进入门厅之后，她按了一下写着杰弗瑞·盖伯名字的对讲机说：“披萨15分钟到。”

杰弗瑞住在柔伊楼上，是一位自由工作者，帮忙研发社群媒体的软件APP。他也提供一些技术咨询来帮忙负担租金，比方如何加快搜索速度、如何在脸书上打广告等。当然，他的发财指望就是日后可以发明一款当红APP。这些年来，他常常提议柔伊来合作一些项目，每一个他提议的项目，好像都会成为下一个Instagram似的。不过，柔伊一直都没有合作意愿。

迄今，他们谁都没有成为下一个Instagram的开发人。

杰弗瑞为人不错，柔伊喜欢他，但受不了他总是愤世嫉俗的样子。对柔伊来说，他对“富人”会自然产生一种敌意，而且特别针对成功富有的企业主。就拿她去上班的企业就是其中一例。尽管她很好奇，如果哪天杰弗瑞的APP也找到一个黄金市场，他不是也会变成这些成功企业家之一吗？尽管如此，他们依然是好朋友，他们发展出每周一起共进一次晚餐的相处模式，轮流作东。

今天晚上轮到柔伊。路易吉家的披萨，经典口味，小小披萨上有各种配料。这是布鲁克林区最好的披萨，一通电话搞定，而且不用洗碗盘。

坦白说，这是柔伊和杰弗瑞少有的共同之处。杰弗瑞不煮东西，柔伊也一样，最多偶尔烤个贝果，或者做个超大的欧姆蛋。柔伊的妈妈也对下厨没兴趣，他们家用的更多的是冷冻柜和微波炉，而不是冷藏室和烹煮炉。柔伊的外婆喜欢烹饪，到了妈妈那里就造反了。“烘焙？我甚至不会做糖霜。”妈妈总爱这么说。

他们一起享用披萨时，柔伊把咖啡师的古怪论点告诉了杰弗瑞。

亨利身上有种特别的气质，吸引着柔伊，简直可说是磁力了。也可以说魅力，但又有奇怪的原因，她无法靠近。就像那张非常吸引她的照片一样。

杰弗瑞听着柔伊的叙述，一声也不吭。

柔伊把她那份吃完后，擦干净手指，拿出笔电。白天上班时，她已经抽空重绘了亨利画的表格，也就是理论上说，如果她1天存25美元，1年就可以存到将近6,800美元，40年后，更是一个天文数字的存款。

她从电脑中调出那张表格来，分享给杰弗瑞。

“杰弗瑞，你看，”她说，“40年后，我67岁，也就是法定退休年纪的时候，超过300万美元的存款耶！”

她的好友也擦干净手指，然后往后靠了靠，再舔干手指，看着柔伊。

“真的假的？”他说：“柔伊，给我一点时间。10\%？怎么可能？哪里可以有10\%的利息？那是多久之前的事了。”

“此外，整个体系也有问题。柔伊，你知道的。你愈是想存钱，政府就拿走了愈多。”

可是，亨利说了，要在政府课税之前，就先存起来一些？对此，柔伊尚未弄懂。

“还有一个，就是通货膨胀。你知道，40年后100万值多少吗？可以和老房东租一个床位就不错了。还有401退休金计划，其中诸多的限制和条款，都是坑我们的。”他又补充说：“还有，谁知道你会在这家公司工作多久？如果有一天离开，退休金又不知会发生什么变化？”

“我无意冒犯，但柔伊，他到底多有知识？那家伙几岁了？60？70？还在煮咖啡？”

对于这些问题，柔伊也没有答案。

杰弗瑞谢过柔伊的披萨，回去楼上。柔伊用了45分钟洗头，清理几乎没什么存货的冰箱，洗刷了一下原本就不太使用的炉台。直到她累得把自己塞到堆满东西的沙发上，她才明白什么叫洁癖。

她不想一直绕在朋友的怀疑论调中。

亨利画的图表，关于“每个人来到地球上，都有特别的使命”的说法，都依然在激励着她。柔伊必须承认，亨利的话激起了她人生的希望。

杰弗瑞却把这些希望都击毁了。

像那样的复利不过是老掉牙的东西。

她拿起手机，在最常联系人当中点击了老妈，电话响了4、5声后，老妈接了电话。

“嗨，老妈。”

“嗨，宝贝女儿。都好吗？”老妈的声音听起来很没力。

“我都好。老妈，你那边呢？你的声音听起来好疲惫。”

“就是讨厌的流感。”老妈回答：“先是你老爸中了，后来，我猜想，流感一定觉得我更可爱一点。”老妈叹了一口气，但是她再说话时，柔伊听得出她声音中的笑意。“听到你的电话，我就已经觉得好多了。你最近都好吗？工作上也都顺利吗？”

“都好，没问题的。我可以请教一个问题吗？老爸公司帮他办了401退休帐户？如果他换工作，会怎么样吗？”

“噢，柔伊，我还真不知道。”老妈说：“那都是你老爸自己在处理的。你该不是想要换工作吧？”

柔伊听到有其他电话插播进来的声音。她看了看屏幕，是洁西卡。“现在我还不知道，妈妈，对不起，我要先挂电话了。”

“宝贝，为已经拥有的而喜乐。远山上的草没有比较绿。”

“好了，我知道了。我明天会再打给你的。我现在要挂电话了，爱你喔！”

说完，柔伊就挂上了妈妈的电话，但就在她要接起洁西卡电话时，她突然犹豫了一下。她也说不上什么原因，但是现在她还不想和洁西卡通话。她让电话自动转去了语音信箱。

她看到新留言的灯在闪，她才拿起话筒，播放留言。

“哈啰，柔伊。我们约周五，对吗？这周该我请客。我们就周五见啰！对了，你和老板谈过要离开杂志社的事吗？”

柔伊挂上电话，重新设置好。“没有，”她对着空荡荡的公寓说：“还没有。”

柔伊再次考虑了经纪人工作的事，其实这个礼拜已经是她本周的第100次考虑了，这个高薪水职位后面的高风险，高压力。她深深吸了一口气，又缓缓吐出来。

洁西卡是生活步调快的人，这是无庸置疑的。如果杰弗瑞职场哲学是忙中偷闲，那么，洁西卡的职场哲学则是把饼做得更大一点。在职场上，她根本不是一级一级往上走，而是一路冲到顶，然后高高站在云端上。

那么柔伊，她自己秉持的职场哲学是哪个呢？
